\section{Introduction}
\label{sec:intro}

Humans interact with the physical world as active agents: we move our viewpoint, manipulate objects, and reason about how actions lead to consequences. Yet existing video generation models lack this unified capability~\cite{wan2025wan, CogVideo, AlignYourLatents, CogVideoX, VideoCrafter2}. Bridging this gap is essential for applications that demand interactive visual reasoning, from embodied AI agents~\cite{bar2025navigation, yang2023learning, wang2024drivedreamer} that must anticipate action outcomes~\cite{hafner2019dream, hafner2023mastering, gao2026dreamdojo}, to world models~\cite{hong2025relic, Sora, he2025matrix, yang2025matrix, ye2026world} that simulate physical interactions, and to immersive content creation~\cite{liu2025ponimator, liu2024physgen, chen2025physgen3d, li2025wonderplay, gao2024gaussianflow, wu2024draganything, zhang2024physdreamer} where users freely navigate and manipulate scenes. A desirable video generation system should therefore offer joint controllability over both camera and object motion---generating visually coherent frames under arbitrary viewpoint changes while producing causally consistent scene dynamics driven by user-specified actions.

Existing controllable video generation methods~\cite{MotionCtrl, geng2024motionprompting, wang2025ati, chu2025wanmove, burgert2025go, wang2024boximator} work as {\it renderers}: given displacements for {\it all} pixels, they generate a visually realistic video that adheres to the displacements.
These approaches have two key limitations in practice. First, they entangle camera and object motion, making joint control ambiguous because viewpoint changes alter pixel trajectories. Second, they interpret user-specified motion as simple kinematic displacement and largely ignore the consequences of the given trajectories. Models, therefore, focus on following trajectories rather than reasoning about causal relationships between objects. In reality, actions cause consequences—pushing a cup may cause it to slide and collide with other objects, while lifting a teapot may cause water to pour. Specifying the effects of all the objects' motion through motion prompts is often impractical. Without modeling these action–consequence relationships, motion-controlled generation cannot capture the causal structure of real-world interactions.

To overcome these challenges, we introduce \ours, a unified framework for video generation with disentangled camera–object motion control and motion causality reasoning as shown in \cref{fig:teaser}.
Given a reference image, user-specified motion trajectories, and target viewpoints, \ours generates videos where objects follow the desired motion and the scene is rendered from the specified cameras.
For disentangled camera-object motion control, our key insight is that specifying the motion of objects under camera changes is inherently difficult. We therefore introduce a \textbf{dual-stream motion formulation}. The first branch models and generates object motion in the source image plane under a canonical static viewpoint, allowing users to easily specify dynamic trajectories. The second branch represents the target camera motion and transfers object dynamics from the canonical branch via temporal cross-view attention. This \emph{cross-view motion transfer} enables independent control of camera and object motion while maintaining coherent scene dynamics.
For motion causality reasoning,
we achieve this by decomposing object motion into two categories during training: \emph{active motion}, representing user-driven actions, and \emph{passive motion}, representing their causal outcomes. By conditioning on either active or passive motion signals, the video model learns to generate all the scene dynamics, capturing action--consequence relationships within the scene. This yields two complementary capabilities: forward reasoning of scene evolution from user actions, and inverse reasoning of plausible actions that produce a desired outcome.

We evaluate \ours on three benchmarks covering diverse interaction scenarios. Results show that \ours outperforms existing methods in generation quality, motion controllability, and interaction awareness, validating the effectiveness of disentangled motion control and causal motion reasoning.

In summary, our contributions are threefold.
(1)~We propose a disentangled framework for camera and object motion control, enabling users to draw motion trajectories directly in the image plane while freely adjusting viewpoints to generate coherent videos.
(2)~We empower video generation models with motion reasoning capability by modeling action--consequence relationships, allowing user-driven motions to meaningfully interact with the environment and produce consistent scene dynamics.
(3)~We demonstrate that \ours supports both forward and inverse reasoning: given active motion inputs, it predicts future scene evolution; given desired passive outcomes, it recovers plausible actions that achieve them. Together, these advances establish a unified framework for controllable and reasoning-aware video generation.
